\documentclass[
a4paper,
12pt,
parskip=half,
numbers=noendperiod,
DIV=15,% ! wenn verändert komplette Arbeit überprüfen!
% BCOR=6mm,% Bindekorrektur
% areasetadvanced,% Bindekorrektur verwenden ohne typearea zu verkleinern
% headinclude,
% headheight=36pt,
toc=chapterentrywithdots
]{scrbook}
% \areaset[12mm]{\textwidth}{\textheight} % s. o., typearea für Bindekorrektur verschieben

%%% Präambel ------------------------------------------+
%%% common packages                                    |
%%% Präambel #######################################################
\usepackage[utf8]{inputenc}                                         % Anpassung der Eingabekodierung
\usepackage[T1]{fontenc}                                            % Standard-Paket für die Auswahl der Schriftarten Kodierung
\usepackage{scrhack}                                                % Paket um Warnungen/Fehler durch veraltete Funktionen bei der Verwendung von KOMA Skript zu beheben (hier: Warnung bei erstem frontmatter Kapitel nach Titelseite)
\usepackage[automark,headsepline]{scrlayer-scrpage}                 % weitere Anpassungen der Dokumentenklasse
\usepackage[onehalfspacing]{setspace}                               % Zeilenabstand
\usepackage[english,ngerman]{babel}                                 % Sprachpaket ("english" ist in diesem Fall nur für den Text im "Abstract" notwendig, kann also ggf. entfernt werden; letzte Sprache hat Priorität)
\usepackage{csquotes}                                               % verbesserter Syntax und Verhalten von Anführungszeichen (empfohlen für biblatex)
\usepackage[style=ieee,sorting=none,giveninits=true,backend=biber]{biblatex} % Literaturverzeichnis und -verweise
\DeclareFieldFormat{title}{\textsl{#1}}                             % Titel im Literaturverzeichnis in schräger statt kursiver Schrift
\DeclareFieldFormat{booktitle}{\textsl{#1}}                         % s. o.
\DeclareFieldFormat{booksubtitle}{\textsl{#1}}                      % s. o.
\DeclareNameAlias{author}{family-given}                             % Nachnamen vor Vornamen anzeigen (vorher: last-first)
\DeclareNameAlias{editor}{family-given}                             % s. o.
\DeclareNameAlias{translator}{family-given}                         % s. o.
\DeclareNameAlias{default}{family-given}                            % s. o.
% \setcounter{biburllcpenalty}{7000}                                  % work-around für unschön umgebrochene URLs im Literaturverzeichnis
% \setcounter{biburlucpenalty}{7000}                                  % s. o.
% \setcounter{biburlnumpenalty}{7000}                                 % s. o.
\usepackage{etoolbox}                                               % Erstellung eigener Dokumentoptionen
\usepackage{lmodern}                                                % Schriftartenpaket
\usepackage{amsmath,amsfonts,amssymb}                               % AMS Pakete
\numberwithin{figure}{chapter}                                      % Kapitelnummer in Abbildungsbeschriftung
\numberwithin{table}{chapter}                                       % Kapitelnummer in Tabellenbeschriftung
\numberwithin{equation}{chapter}                                    % Kapitelnummer in Formelbeschriftung
\usepackage{siunitx}                                                % automatische und korrekte Schreibweise der SI Einheiten
\sisetup{
    locale=DE,                                                      % setzt \cdot statt \times in Exponentialschreibweise und output-decimal-marker = {,}
    group-digits = integer,
    group-separator = {.},                                          % Punkt als 1000er-Trennzeichen
    group-minimum-digits = 4,                                       % schon ab 4 Stellen gruppieren
    per-mode = symbol-or-fraction,                                  % per-mode wird in Abhängigkeit der Umgebung gesetzt, bei Fehlern: [per-mode=fraction] manuell setzen oder bei Drehzahlen: [per-mode=reciprocal]
    % product-symbol = \cdot,                                         % Zeichen bei Produkten, default: \times (x)
    list-final-separator = { und },
    list-pair-separator = { und },
    % list-units = single,                                          % Einheit nur beim letzten Listenelement; kann zu Missverständnissen führen
    % range-units = single,                                         % Einheit nur bei der größeren Bereichsgrenze; kann zu Missverständnissen führen
    range-phrase = { bis }
}
\DeclareSIUnit\dekade{Dekade}
% \DeclareSIPrefix\kilo{k}{3}
% \DeclareSIPower\quartic\tothefourth{4}                              % Bsp.: \unit{\kilogram\tothefourth} oder \unit{\quartic\metre}
% \DeclareSIQualifier\peakpeak{PP}
\usepackage[siunitx,EFvoltages]{circuitikz}                         % Option "EFvoltages" zeigt Spannungspfeile in Richtung des el. magn. Felds (also von + nach -)
\usepackage{graphicx}                                               % Grafiken einbinden
% \graphicspath{{img/}}                                               % Pfad für Grafiken (hier: Ordner "img" im Verzeichnis des Hauptdokuments)
% \usepackage{svg}                                                    % SVG-Grafiken einbinden
% \svgpath{svg/}                                                      % Pfad für SVG-Grafiken (hier: Ordner "svg" im Verzeichnis des Hauptdokuments)
\usepackage{caption}                                                % Beschriftungen an Abbildungen, Tabellen, etc.
\usepackage{subcaption}                                             % mehrere Abbildungen als eine darstellen
\usepackage{rotating}                                               % gedrehte Grafiken und Tabellen
\usepackage{xcolor}                                                 % Farbdefinitionen
\definecolor{MATLABplot1}{HTML}{0072BD}                             % hier beispielhaft aktuelle Standardfarben in MATLAB Plots
\definecolor{MATLABplot2}{HTML}{D95319}                             % https://de.mathworks.com/help/matlab/ref/plot.html#btzitot-Color
\definecolor{MATLABplot3}{HTML}{EDB120}
\definecolor{MATLABplot4}{HTML}{7E2F8E}
\definecolor{MATLABplot5}{HTML}{77AC30}
\definecolor{MATLABplot6}{HTML}{4DBEEE}
\definecolor{MATLABplot7}{HTML}{A2142F}
\usepackage{booktabs}                                               % Tabellen in Buchdruckqualität (horizontale Linien)
% \usepackage{multirow}                                               % mehrzeilige Zellen
% \usepackage{colortbl}                                               % farbige Zellen
\usepackage{listings}                                               % Code-Listings mit automatischem Syntax-Highlighting
\definecolor{codegray}{rgb}{0.5,0.5,0.5}
% \definecolor{MATLABtext}{HTML}{000000}
\definecolor{MATLABbackground}{HTML}{FFFFFF}
\definecolor{MATLABkeywords}{HTML}{0000FF}
\definecolor{MATLABstrings}{HTML}{A020F0}
\definecolor{MATLABcomments}{HTML}{3C763D}
\lstdefinestyle{MATLAB}{
    backgroundcolor=\color{MATLABbackground},
    commentstyle=\color{MATLABcomments},
    keywordstyle=\color{MATLABkeywords},
    stringstyle=\color{MATLABstrings},
    % identifierstyle=\color{blue},% hier wird auch das e in bspw. 10e-3 erkannt
    numberstyle=\tiny\color{codegray},
    basicstyle=\ttfamily\footnotesize,
    breakatwhitespace=false,
    breaklines=true,
    captionpos=b,
    keepspaces=true,
    numbers=left,
    numbersep=5pt,
    showspaces=false,
    showstringspaces=false,
    showtabs=false,
    tabsize=4,
    frame=single
}
\lstset{literate=%
    {Ö}{{\"O}}1
    {Ä}{{\"A}}1
    {Ü}{{\"U}}1
    {ß}{{\ss}}1
    {ü}{{\"u}}1
    {ä}{{\"a}}1
    {ö}{{\"o}}1
}
% \lstset{literate={~}{{\textasciitilde}}1} % wenn benötigt, in obere Liste einfügen
% \lstset{literate={µ}{{$\mathrm{\mu}$}}1} % s. o.

% ^^========= weitere Pakete über dieser Zeile einfügen ==========^^
% * Paket hyperref (und bookmark) immer als letztes laden! Paket glossaries nach hyperref laden!
\usepackage[ngerman]{hyperref}                                      % Hyperlinks in PDF Datei
\usepackage{bookmark}                                               % Lesezeichen in PDF Datei
% "You must load the hyperref package before the glossaries package to ensure the hyperlinks work." [S. 124, glossaries-user.pdf]
\usepackage[acronym,nopostdot,nonumberlist,toc]{glossaries}         % Glossar und Abkürzungen; für Eintrag im Inhaltsverzeichnis Option "toc" hinzufügen; "automake" "nomain" ?
% Änderung des Darstellungsformats und Berechnung des breitesten Eintrags: siehe "User Manual for glossaries.sty v4.46", Seite 222
\renewcommand{\glossarypreamble}{%
    \glsfindwidesttoplevelname[\currentglossary]}
\setglossarystyle{alttree}
% Code für Änderung des Darstellungsformats und Berechnung des breitesten Eintrags bis hier
\makeglossaries
% oder:
% \usepackage[nopostdot,nonumberlist]{glossaries}
% \usepackage[nohyperlinks, printonlyused]{acronym}
% Formatierung: (Bsp., weitere unter: https://www.dickimaw-books.com/gallery/glossaries-styles/)
% \usepackage{glossary-longragged}
% \setglossarystyle{longragged}
% https://tex.stackexchange.com/questions/415275/flush-glossaries-super-style-to-left

% Anpassung der Kopf- und Fußzeilen ================================
% \clearpairofpagestyles
% \lohead{...}
% \rohead{...}
% \cfoot[...]{...}
% \pagestyle{scrheadings}

% Anpassung der automatischen Bezeichnungen ========================
\renewcaptionname{ngerman}\sectionautorefname{Abschnitt}
\renewcaptionname{ngerman}\subsectionautorefname{Abschnitt}
\renewcaptionname{ngerman}\subsubsectionautorefname{Abschnitt}

% Anpassung der Seitenränder auf der Titelseite (funktioniert nur für \maketitle und nicht titlepage Umgebung)
% \renewcommand*{\coverpagetopmargin}{0mm}
% \renewcommand*{\coverpageleftmargin}{0mm}
% \renewcommand*{\coverpagerightmargin}{0mm}
% \renewcommand*{\coverpagebottommargin}{0mm}

% Anpassung der PDF Eigenschaften ==================================
\AtEndPreamble{
    \hypersetup{
        pdftitle    = {\VarTitel},
        pdfsubject  = {\VarArtDerArbeit{} im Studiengang \VarStudiengang},
        pdfauthor   = {\VarVerfasserIn},
        % pdfkeywords = {Stichwort1, Stichwort2 ...},
        % pdfdisplaydoctitle = true,                                      % Dokumenttitel statt Dateiname im PDF Viewer anzeigen
        % pdfpagelayout=TwoPageRight,                                     % Voreinstellung Ansicht: "TwoColumnRight" kontinuierliches Scrollen, "TwoPageRight" seitenweises Scrollen
        bookmarksnumbered = true,                                       % Abschnittsnummerierung in Lesezeichen einbeziehen
        bookmarksopen = true,                                           % Lesezeichen aufgeklappt anzeigen
        % colorlinks = true,                                              % Links farbig darstellen
        hidelinks,                                                      % Textfarbe von Links: schwarz
        unicode=true
    }
}

% neue Befehle für Dokumentfelder ==================================
\def\VarLogoHsH{7}
\newcommand{\LogoHsH}[1]{
    \def\VarLogoHsH{#1}
}
\newcommand{\Fakultaet}[1]{
    \def\VarFakultaet{#1}
}
\newcommand{\Fachgebiet}[1]{
    \def\VarFachgebiet{#1}
}
\newcommand{\ArtDerArbeit}[1]{
    \def\VarArtDerArbeit{#1}
}
\newcommand{\Studiengang}[1]{
    \def\VarStudiengang{#1}
}
\newcommand{\Titel}[1]{
    \def\VarTitel{#1}
}
\newcommand{\Untertitel}[1]{
    \def\VarUntertitel{#1}
}
\newcommand{\VerfasserIn}[3][Verfasser:in]{
    \def\VarBezVerfasserIn{#1}%
    \def\VarVerfasserIn{#2}%
    \def\VarVerfasserInMatNr{#3}
}
\newcommand{\DatumBeginn}[2][Beginn]{
    \def\VarBezDatumBeginn{#1}%
    \def\VarDatumBeginn{#2}
}
\newcommand{\DatumAbgabe}[2][Abgabe]{
    \def\VarBezDatumAbgabe{#1}%
    \def\VarDatumAbgabe{#2}
}
\newtoggle{DatumLeerzeile}
\newcommand{\DatumLeerzeile}{
    \toggletrue{DatumLeerzeile}
}
\newtoggle{DatenAlsZeitraum}
\newcommand{\DatenAlsZeitraum}[1][Zeitraum]{
    \def\VarBezDatenZeitraum{#1}%
    \toggletrue{DatenAlsZeitraum}
}
\newcommand{\BetreuerInA}[2][Betreuer:in]{
    \def\VarBetreuerInA{#2}%
    \def\VarBezBetreuerInA{#1}
}
\newcommand{\BetreuerInB}[2][]{
    \def\VarBetreuerInB{#2}%
    \def\VarBezBetreuerInB{#1}
}
\newcommand{\BetreuerInC}[2][]{
    \def\VarBetreuerInC{#2}%
    \def\VarBezBetreuerInC{#1}
}
\newcommand{\DatumVersion}[1]{
    \def\VarDatumVersion{#1}
}
\newtoggle{BetreuerInLeerzeileAB}
\newcommand{\BetreuerInLeerzeileAB}{
    \toggletrue{BetreuerInLeerzeileAB}
}
\newcommand{\BetreuerInLeerzeileBC}{
    \toggletrue{BetreuerInLeerzeileBC}
}


% [alle] Indizes Steil Stellen (von Cube, HsH-Forum)
% https://betech-hsh.de/viewtopic.php?f=119&t=756
\def\subinrm#1{\sb{\mathrm{#1}}}
{\catcode`\_=13 \global\let_=\subinrm}
\mathcode`_="8000
% Macros definieren:
\def\upsubscripts{\catcode`\_=12 }
\def\normalsubscripts{\catcode`\_=8 }
% Macro benutzen:
\upsubscripts %Indizes Steil

%%% packages for development                           |
%%% Pakete, die nur bei der Dokumenterstellung benötigt werden #####
\usepackage{blindtext}                                              % nur für Blindtexte benötigt, kann bei Bedarf entfernt werden
% \usepackage{showframe}                                              % zeigt die Textbereiche (und Seitenränder) an
% \usepackage[mark]{gitinfo2}
%%% packages with draft settings                       |
% \input{tex/prae-draft}
%%% packages with final settings/for final compilation |
% \input{tex/prae-final}
%%% ---------------------------------------------------+

%%% manuelle Anpassungen der Worttrennungen
%%% manuelle Anpassungen der Worttrennungen
\hyphenation{
Feld-ef-fekt-tran-sis-tor % Beispiel für Worttrennungen an definierten Stellen
MOSFET % Beispiel um automatische Worttrennung zu unterdrücken
}

%%% Glossar bzw. Akronyme einbinden
% Beispiel
\newglossaryentry{moodle}
{
        name=Moodle,
        description={\acrfull{lms}, das verschiedene Werkzeuge zur Organisation und Durchführung von Lehrveranstaltungen bietet und somit zeit- und ortsunabhängiges Lernen in virtuellen Kursräumen ermöglicht \cite{HsH.ELC}} % verschachtelte Glossareinträge vermeiden, Akronyme können jedoch verwendet werden
}
% Beispiele
\newacronym{hsh}{HsH}{Hochschule Hannover}
\newacronym[longplural={Lernmanagementsysteme}]{lms}{LMS}{Lernmanagementsystem}
\newacronym[longplural={elektronische Pedalwertgeber}]{epwg}{ePWG}{elektronischer Pedalwertgeber}


%%% Literatur-Liste einbinden - aus Kompatibilitätsgründen mit LaTeX Workshop (VS Code Erweiterung) hier
\addbibresource{bib/Literatur.bib}

%%% hier Titel, Namen, Daten definieren --------------------+
%%% nachfolgenden Kommentar entfernen um HsH Logo auf       |
%%% Titelseite einzufügen; Zahl im Argument des Befehls     |
%%% entspricht Fakultät: 0 -> grau, 1 -> Fak. I, usw.       |
\LogoHsH{1}
\Fakultaet{Fakultät I -- Elektro- und Informationstechnik}
\Fachgebiet{Fachgebiet Regelungstechnik und Mechatronik}
\ArtDerArbeit{Masterarbeit}
\Studiengang{Elektrische Energiesysteme und Elektromobilität}
\Titel{Reglerentwurf und Simulation für die Geschwindigkeitsregelanlage eines Elektrokleinfahrzeugs}
\Untertitel{im Rahmen des Forschungsprojekts XY}
% \Untertitel{im Zusammenarbeit mit\\[4pt]\includegraphics[width=0.1\textwidth]{example-image}}
\VerfasserIn[Verfasserin]{Jane Doe}{1234567}
%%% über optionales Argument kann auch hier die Bezeichnung |
%%% des Feldes auf der Titelseite angepasst werden          |
\DatumBeginn[Bearbeitungsbeginn]{1. September 2023}
%%% über optionales Argument kann auch hier die Bezeichnung |
%%% des Feldes auf der Titelseite angepasst werden          |
\DatumAbgabe[Abgabe der Arbeit]{29. Februar 2024}
%%% Kommentar entfernen um Leerzeile zwischen Daten         |
%%% einzufügen                                              |
% \DatumLeerzeile
%%% Kommentar entfernen um Daten als Zeitraum anzugeben;    |
%%% über optionales Argument kann auch hier die Bezeichnung |
%%% des Feldes auf der Titelseite angepasst werden          |
\DatenAlsZeitraum[Bearbeitungszeitraum]
\BetreuerInA[Erstprüferin]{Prof. Dr.-Ing. Erika Musterfrau}
\BetreuerInB[Zweitprüfer]{Prof. Dr.-Ing. Max Mustermann}
\BetreuerInC[Firmenbetreuerin]{Frieda Musterfrau, M.Sc.}
%%% Kommentar entfernen um Leerzeile zwischen Betreuer:innen|
%%% A und B einzufügen                                      |
% \BetreuerInLeerzeileAB
%%% Kommentar entfernen um Leerzeile zwischen Betreuer:innen|
%%% B und C einzufügen                                      |
% \BetreuerInLeerzeileBC
%%% Datumsangabe in Kopfzeile, mit "\VarDatumAbgabe" statt  |
%%% "\today" wird das Datum verwendet, welches mit          |
%%% "\DatumAbgabe{}" definiert wurde                        |
\DatumVersion{\today}
%%% --------------------------------------------------------+

%%% nur hier angegebene Kapitel werden dargestellt, aber Metadaten aller anderen werden mit einbezogen
%%% ! Literaturverzeichnis zeigt nur Quellen der angegeben Kapitel --> Nummerierung nicht final, später nochmals überprüfen!
% \includeonly{1-Einleitung,3-Schluss}

\begin{document}

    %%% Titelblatt (Cover) ================================ % fehlt dann bei gebundener Ausgabe

    %%% Schmutztitel ======================================
    % \extratitle{Schmutztitel}

    %%% Titelrückseite ====================================
    % \frontispiece{Frontispiz}

    %%% Vorwort ###########################################
    \frontmatter
    \pagenumbering{Roman}

    %%% Titelblatt (keine Bilder, etc.) ===================
    \KOMAoptions{twoside = false}
\begin{titlepage}
    \begin{center}
        \ifnumequal{\VarLogoHsH}{0}{\vspace*{-48pt}\includegraphics[width=0.2\textwidth]{hsh/HsH-Logo-grau}\\[32pt]}{}
        \ifnumequal{\VarLogoHsH}{1}{\vspace*{-48pt}\includegraphics[width=0.2\textwidth]{hsh/HsH-Logo-F1}\\[32pt]}{}
        \ifnumequal{\VarLogoHsH}{2}{\vspace*{-48pt}\includegraphics[width=0.2\textwidth]{hsh/HsH-Logo-F2}\\[32pt]}{}
        \ifnumequal{\VarLogoHsH}{3}{\vspace*{-48pt}\includegraphics[width=0.2\textwidth]{hsh/HsH-Logo-F3}\\[32pt]}{}
        \ifnumequal{\VarLogoHsH}{4}{\vspace*{-48pt}\includegraphics[width=0.2\textwidth]{hsh/HsH-Logo-F4}\\[32pt]}{}
        \ifnumequal{\VarLogoHsH}{5}{\vspace*{-48pt}\includegraphics[width=0.2\textwidth]{hsh/HsH-Logo-F5}\\[32pt]}{}
        \ifnumequal{\VarLogoHsH}{6}{\vspace*{-48pt}\includegraphics[width=0.2\textwidth]{hsh/HsH-Logo-gelb}\\[32pt]}{}
        \textbf{\LARGE{Hochschule Hannover}}\\
        \Large{University of Applied Sciences and Arts}\\[16pt]
        \textbf{\large{\VarFakultaet}}\\
        \large{\VarFachgebiet}\\[32pt]
        \textbf{\Large{\VarArtDerArbeit}}\\
        \large{im Studiengang \VarStudiengang}\\[20pt]
        \textbf{\LARGE{\VarTitel}}\\[12pt]
        \large{\VarUntertitel}\\[2pt]
    \end{center}
    
    \begin{flushleft}
        \begin{table}[b!]
            \begin{tabular}{lll}
                \textbf{\VarBezVerfasserIn} & & \VarVerfasserIn, \VarVerfasserInMatNr \\
                & & \\
                \iftoggle{DatenAlsZeitraum}{
                    \textbf{\VarBezDatenZeitraum} & & \VarDatumBeginn{} -- \VarDatumAbgabe \\ % true
                }{
                    \textbf{\VarBezDatumBeginn} & & \VarDatumBeginn \\
                    \iftoggle{DatumLeerzeile}{
                        & & \\
                    }{
                        %
                    }
                    \textbf{\VarBezDatumAbgabe} & & \VarDatumAbgabe \\ % false
                }
                & & \\
                \textbf{\VarBezBetreuerInA} & & \VarBetreuerInA \\
                \iftoggle{BetreuerInLeerzeileAB}{
                    & & \\
                }{
                    %
                }
                \textbf{\VarBezBetreuerInB} & & \VarBetreuerInB \\
                \textbf{\VarBezBetreuerInC} & & \VarBetreuerInC \\
            \end{tabular}
        \end{table}
    \end{flushleft}
\end{titlepage}
\KOMAoptions{twoside}
    \cleardoublepage

    %%% Erklärung zur ordnungsgemäßen Erstellung ==========
    %%% alternativ: gedrucktes Formular der HsH nutzen und beifügen
    %%% Version HsH
    \thispagestyle{plain.scrheadings}
    \section*{Eigenhändigkeitserklärung}
    Ich versichere, dass ich meine am \VarDatumAbgabe{} vorgelegte \VarArtDerArbeit{} selbstständig verfasst und alle benutzten Hilfsmittel angegeben habe.\\[48pt]
    Hannover, den \VarDatumAbgabe\\[40pt]
    \begin{tabular}{@{}p{5cm}l}
        \hline
        \VarVerfasserIn, \VarVerfasserInMatNr
    \end{tabular}
    \clearpage

    %%% Danksagung (oder: Vorwort und Dank) ===============
    %%% Vorwort und Dank [preface and acknowledgement]
\chapter{Vorwort und Dank}
\blindtext % erzeugt Blindtext, kann entfernt werden % erzeugt Blindtext, kann entfernt werden

    \cleardoublepage

    %%% Kurzfassung/Abstract ==============================
    %%% Kurzfassung [abstract]
\chapter{Kurzfassung}
\blindtext % erzeugt Blindtext, kann entfernt werden

\subsection*{Schlagworte}
Blindtext, Test, Textausgabe


\begin{otherlanguage}{english}

\section*{Abstract}
\blindtext % erzeugt Blindtext, kann entfernt werden

\subsection*{Keywords}
dummy text, test, text output

\end{otherlanguage}

    \cleardoublepage

    %%% Abkürzungen und Glossar (mit Paket glossaries) ====
    \thispagestyle{plain.scrheadings}
\phantomsection
\printacronyms[title=Abkürzungen] % \printglossary[type=\acronymtype,title=Abkürzungen]   % alternativ: \printacronyms[title=Abkürzungen]; optional: "toctitle=..."

%%% ggf. stattdessen manuelles Abkürzungsverzeichnis
% \phantomsection
% \input{Dateiname}

\phantomsection
\printglossary[title=Glossar]
    \cleardoublepage

    %%% Symbolverzeichnis =================================
    %%% Symbolverzeichnis (Formelzeichen und Indizes)
\chapter{Symbolverzeichnis}

\section*{Formelzeichen}
% Table generated by Excel2LaTeX from sheet 'Formelzeichen' from file 'SymbolVZ'
\begin{table}[htbp]
    \centering
        \begin{tabular}{lll}
            \toprule
            Formelzeichen   & Name                              & SI-Einheit \\
            \midrule
            $I$             & elektrische Stromstärke           & \si{\ampere} \\
            $R$             & elektrischer Widerstand           & \si{\ohm} \\
            $U$             & elektrische Spannung              & \si{\volt} \\
            $\omega$        & Kreisfrequenz                     & \si{\radian\per\second} \\
            \bottomrule
        \end{tabular}
\end{table}

\clearpage

\section*{Indizes}
% Table generated by Excel2LaTeX from sheet 'Indizes' from file 'Symbol_VZ'
\begin{table}[htbp]
    \centering
        \begin{tabular}{lll}
            \toprule
            Index                  & Bedeutung              & Übersetzung (ggf.) \\
            \midrule
            in                     & input                  & Eingang \\
            max                    & maximum                & Maximum \\
            min                    & minimum                & Minimum \\
            out                    & output                 & Ausgang \\
            \bottomrule
        \end{tabular}
\end{table}
    \cleardoublepage

    %%% Inhaltsverzeichnis ================================
    \pagestyle{plain.scrheadings}
    \phantomsection
    % \addcontentsline{toc}{chapter}{Inhaltsverzeichnis}
    \tableofcontents
    \cleardoublepage

    %%% Hauptteil #########################################
    \mainmatter
    \pagestyle{scrheadings}
    \pagenumbering{arabic}
    \setcounter{page}{1}

    \chapter{Einleitung}
\label{chap:Einleitung}

\blindmathpaper % erzeugt Blindtext mit mathematischen Formeln, kann entfernt werden % erzeugt Blindtext mit mathematischen Formeln, kann entfernt werden
\blindmathpaper % erzeugt Blindtext mit mathematischen Formeln, kann entfernt werden


    \section{Hauptteil}
\label{sec:Hauptteil}

\subsection{Beispiele für Zitate, Formeln, Tabellen und Abbildungen}

Erstmals wurde der Zusammenhang von elektrischer Spannung, elektrischem Strom und elektrischem Widerstand von Georg Simon Ohm beschrieben. \cite{Ohm.1827}

\begin{equation}
    U = R \cdot I
\end{equation}

\begin{table}[htbp]
    \centering
    \caption{Beispieltabelle}
    \label{tab:bsp-tabelle}
    \begin{tabular}{lll}
        \toprule
        \textbf{Formelzeichen} & \textbf{Name} & \textbf{SI-Einheit} \\
        \midrule
        $U$ & Elektrische Spannung & \si{\volt} \\
        $R$ & Elektrischer Widerstand & \si{\ohm} \\
        $I$ & Elektrischer Strom & \si{\ampere} \\
        \bottomrule
    \end{tabular}
\end{table}

\begin{figure}[htbp]
    \centering
    \includegraphics[width=0.8\textwidth]{example-image} % dieses Bild ist Teil des mwe Pakets, welches auch ohne Befehl in der Präambel aufgerufen wird
    \caption{Beispielbild}
    \label{fig:bsp-bild}
\end{figure}

\subsection{Beispiele für die Nutzung von Abkürzungen und Glossareinträgen}

\begin{table}[htbp]
    \centering
    \caption{Befehle und Ausgaben für Paket \texttt{glossaries}}
    \label{tab:befehle-ausgaben-glossaries}
    \begin{tabular}{lll}
        \toprule
        Befehl & Kurzbefehl & Ausgabe \\
        \midrule
        \verb|\acrfull{epwg}| & \verb|\acf{epwg}| & \acrfull{epwg} \\
        \verb|\acrlong{epwg}| & \verb|\acl{epwg}| & \acrlong{epwg} \\
        \verb|\acrshort{epwg}| & \verb|\acs{epwg}| & \acrshort{epwg} \\
        \verb|\acrfullpl{epwg}| & \verb|\acfp{epwg}| & \acrfullpl{epwg} \\
        \verb|\acrlongpl{epwg}| & \verb|\aclp{epwg}| & \acrlongpl{epwg} \\
        \verb|\acrshortpl{epwg}| & \verb|\acspl{epwg}| & \acrshortpl{epwg} \\
        \verb|\Acrfull{epwg}| & \verb|\Acf{epwg}| & \Acrfull{epwg} \\
        \verb|\Acrlong{epwg}| & \verb|\Acl{epwg}| & \Acrlong{epwg} \\
        \verb|\Acrshort{epwg}| & \verb|\Acs{epwg}| & \Acrshort{epwg} \\
        \verb|\Acrfullpl{epwg}| & \verb|\Acfp{epwg}| & \Acrfullpl{epwg} \\
        \verb|\Acrlongpl{epwg}| & \verb|\Aclp{epwg}| & \Acrlongpl{epwg} \\
        \verb|\Acrshortpl{epwg}| & \verb|\Acspl{epwg}| & \Acrshortpl{epwg} \\
        \bottomrule
    \end{tabular}
\end{table}

Erste Verwendung von \verb|\gls{lms}| führt zu: \gls{lms}

Nächste Verwendung von \verb|\gls{lms}| führt zu: \gls{lms}

Wird ein im Glossar definierter Begriff mit bspw. \verb|\gls{moodle}| verwendet, führt dies zu folgender Ausgabe: \gls{moodle}

\clearpage


    \section{Abschluss}
\label{sec:Abschluss}

\Blindtext % erzeugt viel Blindtext, kann entfernt werden


    \cleardoublepage

    %%% ggf. alle Einträge in Literaturverzeichnis überprüfen
    % \nocite{*}

    %%% Anhang ############################################
    \appendix % geht nur vor \backmatter
    % \section{Anhang}
% \label{sec:Anhang}

\blindmathtrue % aktiviert in-line Formelausdrücke im Blindtext, kann entfernt werden
\blinddocument % Erzeugt ein Dokument mit Blindtext, Überschriften und Listen; kann entfernt werden
\blindmathfalse % deaktiviert in-line Formelausdrücke im Blindtext, kann entfernt werden


    %%% Verzeichnisse #####################################
    \backmatter

    %%% Abbildungsverzeichnis =============================
    \phantomsection
    \addcontentsline{toc}{chapter}{Abbildungsverzeichnis}
    \listoffigures
    \cleardoublepage

    %%% Tabellenverzeichnis ===============================
    \phantomsection
    \addcontentsline{toc}{chapter}{Tabellenverzeichnis}
    \listoftables
    \cleardoublepage

    %%% Literatur-/Quellenverzeichnis =====================
    \phantomsection
    \addcontentsline{toc}{chapter}{Literatur}
    \printbibliography

    %%% Index/Stichwortverzeichnis ========================

\end{document}
